% Part: incompleteness
% Chapter: representability-in-q
% Section: representing-relations

\documentclass[../../../include/open-logic-section]{subfiles}

\begin{document}

\olfileid{inc}{req}{rel}

\olsection{بازنمایی روابط}

بگوییم بازنمایی‌پذیربودن یک \emph{رابطه} به چه معناست.

\begin{defn}
\ollabel{defn:representing-relations} رابطهٔ $R(x_0,\dots,x_k)$ روی
اعداد طبیعی {\em در~$\Th{Q}$ بازنمایی‌پذیر است} هرگاه فرمولی چون
$!A_R(x_0,\dots,x_k)$ وجود داشته باشد به‌گونه‌ای که هرگاه
$R(n_0,\dots,n_k)$ صادق باشد، $\Th{Q}$ فرمول
$!A_R(\num{n_0},\dots,\num{n_k})$ را ثابت کند، و هرگاه
$R(n_0,\dots,n_k)$ کاذب باشد، $\Th{Q}$ فرمول
$\lnot !A_R(\num{n_0}, \dots, \num{n_k})$ را ثابت کند.
\end{defn}

\begin{thm}
\ollabel{thm:representing-rels} یک رابطه در~$\Th{Q}$ بازنمایی‌پذیر
است اگر و تنها اگر محاسبه‌پذیر باشد.
\end{thm}

\begin{proof}
برای جهت رفت، فرض کنید $R(x_0,\dots,x_k)$ با فرمول
$!A_R(x_0,\dots,x_k)$ بازنمایی شود. الگوریتمی برای محاسبهٔ~$R$ چنین
است: با ورودی‌های $n_0$، \dots،~$n_k$، هم‌زمان در جست‌وجوی اثباتی برای
$!A_R(\num{n_0}, \dots, \num{n_k})$ و اثباتی برای
$\lnot !A_R(\num{n_0}, \dots, \num{n_k})$ باشید. بنا بر فرض ما،
جست‌وجو ناگزیر یکی از این دو را می‌یابد؛ اگر اولی را یافت، پاسخ
``بله'' و در غیر این صورت پاسخ ``نه'' بدهید.

در جهت دیگر، فرض کنید $R(x_0, \dots, x_k)$ محاسبه‌پذیر باشد. بنا بر
تعریف، این بدان معناست که تابع $\Char{R}(x_0, \dots, x_k)$
محاسبه‌پذیر است. بنا بر \olref[int]{thm:representable-iff-comp}،
$\Char{R}$ با فرمولی، مثلاً $!A_{\Char{R}}(x_0, \dots, x_k,
y)$، بازنمایی می‌شود. بگذارید $!A_R(x_0, \dots, x_k)$ فرمول
$!A_{\Char{R}}(x_0, \dots, x_k, \num{1})$ باشد. آنگاه برای هر
$n_0$، \dots،~$n_k$، اگر $R(n_0, \dots, n_k)$ صادق باشد،
$\Char{R}(n_0, \dots, n_k) = 1$، که در این حالت $\Th{Q}$ فرمول
$!A_{\Char{R}}(\num{n_0}, \dots, \num{n_k}, \num{1})$ را ثابت
می‌کند، و بنابراین $\Th{Q}$ فرمول $!A_R(\num{n_0}, \dots,
\num{n_k})$ را ثابت می‌کند. از سوی دیگر، اگر
$R(n_0, \dots, n_k)$ کاذب باشد، آنگاه
$\Char{R}(n_0, \dots, n_k) = 0$. این بدان معناست که $\Th{Q}$ ثابت
می‌کند
\[
\lforall[y][(!A_{\Char{R}}(\num{n_0}, \dots, \num{n_k}, y) \lif y =
  \num{0})].
\]
چون $\Th{Q}$ فرمول $\eq/[\num{0}][\num{1}]$ را ثابت می‌کند، $\Th{Q}$ فرمول
$\lnot !A_{\Char{R}}(\num{n_0}, \dots, \num{n_k}, \num{1})$ را نیز
ثابت می‌کند، و بنابراین $\lnot !A_R(\num{n_0}, \dots,
\num{n_k})$ را ثابت می‌کند.
\end{proof}

\begin{prob}
نشان دهید اگر $R$ در~$\Th{Q}$ بازنمایی‌پذیر باشد، آنگاه~$\Char{R}$
نیز در آن بازنمایی‌پذیر است.
\end{prob}

\end{document}
